\chapter{الحركة النسبية}\label{ch:g10:relative-motion}

يتمشى مسافر نحو مقدمة قطار سريع، وفي يده قهوة،
بخطى مشي هادئة. ويسجّل راصد على جانب السكة سرعة ذلك المسافر نفسه
بأكثر من \qty{300}{km/h}. والقياسان صحيحان كلاهما. وليس أيّ منهما
``هي'' سرعة المسافر، إذ لا وجود لشيء كهذا: فالحركة
ليست خاصية للجسم وحده، بل للجسم \emph{و}وجهة
النظر المختارة لوصفه. ويضبط هذا الفصل وجهة النظر تلك
--- ويعلّمك اختيار وجهة جيدة.

\section{الحركة نسبية}

\begin{definition}[المرجع]\label{def:g10:relative-motion:frame}
\emph{المرجع}\index{مرجع} جسم صلب --- الأرض،
أو عربة قطار، أو الكرة الأرضية بأسرها --- تُسجَّل المواضع بالنسبة إليه،
مع ساعة تؤرّخها. ويكون الجسم \emph{ساكنًا}\index{سكون}
في مرجعٍ ما إذا لم يتغير موضعه في ذلك المرجع
مع الزمن؛ وإلا فهو متحرك في ذلك المرجع.
\end{definition}

\begin{definition}[المسار]\label{def:g10:relative-motion:trajectory}
\emph{مسار}\index{مسار} جسمٍ في \omterm{def:g10:relative-motion:frame}{مرجع}
معطًى هو المنحني الذي تشكّله مواضعه المتتابعة في ذلك \omterm{def:g10:relative-motion:frame}{المرجع}:
مستقيم أو دائرة أو قوس\dots{} ولا يكون للمسار معنى
إلا بعد تسمية \omterm{def:g10:relative-motion:frame}{المرجع}.
\end{definition}

\begin{example}[الماشي في القطار]\label{ex:g10:relative-motion:walker}
يسير قطار بسرعة \qty{300}{km/h} بالنسبة إلى الأرض بينما يمشي
مسافر نحو المقدمة بسرعة \qty{4}{km/h} بالنسبة إلى القطار.
\begin{itemize}
  \item في \omterm{def:g10:relative-motion:frame}{مرجع} \emph{القطار}: يتحرك المسافر بسرعة \qty{4}{km/h}
        في الممرّ؛ والمقاعد ساكنة؛ والأشجار في الخارج تندفع
        إلى الوراء بسرعة \qty{300}{km/h}.
  \item في \omterm{def:g10:relative-motion:frame}{المرجع} \emph{الأرضي}: تتحرك المقاعد بسرعة \qty{300}{km/h}،
        ويتحرك المسافر بسرعة \qty{304}{km/h}، والأشجار ساكنة.
\end{itemize}
المشهد نفسه، ومرجعان، ووصفان مختلفان اختلافًا تامًّا --- ولا
تجربة تُجرى داخل القطار تقدر أن تعلن أن أحدهما ``هو
الحقيقي''.
\end{example}

\begin{proposition}[نسبية الحركة]\label{prop:g10:relative-motion:relativity}
حالة السكون أو الحركة لجسمٍ ما، ومساره، وسرعته، كلها
تتوقف على \omterm{def:g10:relative-motion:frame}{المرجع} المختار لوصفها.
\end{proposition}

\begin{proof}
مثال واحد يحسم كل ادّعاء. فالمسافر الجالس في
\cref{ex:g10:relative-motion:walker} ساكن في \omterm{def:g10:relative-motion:frame}{مرجع} القطار
ويتحرك بسرعة \qty{300}{km/h} في \omterm{def:g10:relative-motion:frames}{المرجع الأرضي} (فالسكون والسرعة يتوقفان
على \omterm{def:g10:relative-motion:frame}{المرجع}). أما \omterm{def:g10:relative-motion:trajectory}{المسار}، فراقب صمام عجلة دراجة: يرى
الراصد الراكب للدراجة أنه يرسم دائرة حول المحور، بينما
يرى الراصد الواقف على الطريق أنه يرسم متتابعة من الأقواس (منحنًى
يسمّى \emph{الدويرية}\index{دويرية}) --- انظر الشكل أدناه.
\end{proof}

\begin{omfigure}
\begin{tikzpicture}[scale=0.85, font=\small]
  \draw[gray!60] (-1.6,0) -- (1.6,0);
  \draw[gray, thin] (0,1) circle (1);
  \fill (0,1) circle (1.5pt);
  \draw[gray] (0,1) -- (0.71,1.71);
  \draw[omDef, thick, -Stealth] (1,1) arc (0:330:1);
  \fill[omThm] (0.71,1.71) circle (2pt) node[above right] {الصمام};
  \node[omDef] at (0,-0.5) {مرجع الدراجة: دائرة};
\end{tikzpicture}
\qquad
\begin{tikzpicture}[scale=0.42, font=\small]
  \draw[gray!60] (-0.6,0) -- (13.2,0);
  \draw[omDef, thick] plot[domain=0:12.566, samples=100, smooth]
    ({\x - sin(\x r)}, {1 - cos(\x r)});
  \draw[gray, thin] (9.42,1) circle (1);
  \fill (9.42,1) circle (1.5pt);
  \draw[gray] (9.42,1) -- (9.42,2);
  \fill[omThm] (9.42,2) circle (3pt);
  \node[omDef] at (6.3,-1.0) {مرجع الطريق: دويرية};
\end{tikzpicture}

{\small الصمام نفسه، ومرجعان: دائرة لراكب الدراجة، ودويرية
--- سلسلة أقواس --- للراصد على الطريق.}
\end{omfigure}

\begin{remark}
ولاحظ ما \emph{لا} يتغير في الصورتين: الصمام،
والعجلة، والفيزياء. إنما يتغير الوصف وحده. والسؤال ``هل يدور الصمام
في دائرة حقًّا؟'' كالسؤال عمّا إذا كان جبل
``حقًّا'' إلى اليسار: فذلك يتوقف على موضع وقوفك.
\end{remark}

\section{السرعة المتوسطة}

\begin{definition}[السرعة المتوسطة]\label{def:g10:relative-motion:speed}
\emph{السرعة المتوسطة}\index{سرعة!متوسطة} لجسمٍ ما، في \omterm{def:g10:relative-motion:frame}{مرجع}
معطًى، هي المسافة $d$ التي يقطعها في ذلك \omterm{def:g10:relative-motion:frame}{المرجع} مقسومة على
زمن الانتقال $t$:
\[
v = \frac{d}{t} .
\]
وبأخذ $d$ \omterm{def:g10:orders-of-magnitude:unit}{بالمتر} و$t$ \omterm{def:g10:orders-of-magnitude:unit}{بالثانية}، تكون $v$ \omterm{def:g10:orders-of-magnitude:unit}{بالمتر} في \omterm{def:g10:orders-of-magnitude:unit}{الثانية}
(\unit{m/s})؛ أما الحياة اليومية فتفضّل الكيلومتر في الساعة (\unit{km/h}).
\end{definition}

\begin{method}[التحويل بين m/s وkm/h]\label{met:g10:relative-motion:convert}
الكيلومتر الواحد \qty{1000}{m} والساعة الواحدة \qty{3600}{s}، ومنه
$\qty{1}{m/s} = \qty{3.6}{km/h}$:
\begin{itemize}
  \item من \unit{m/s} إلى \unit{km/h}: \emph{اضرب} في $3.6$؛
  \item من \unit{km/h} إلى \unit{m/s}: \emph{اقسم} على $3.6$.
\end{itemize}
وتحقّق من السلامة: عدد \unit{km/h} هو دائمًا \emph{الأكبر} من
العددين.
\end{method}

\begin{example}[بعض السرعات الجديرة بالمعرفة]\label{ex:g10:relative-motion:convert}
مشي نشيط: $\qty{1.5}{m/s} \times 3.6 = \qty{5.4}{km/h}$. وسيارة مدينة
بسرعة \qty{54}{km/h}: $54 / 3.6 = \qty{15}{m/s}$ --- أي نحو \qty{15}{m}
مقطوعة في كل \omterm{def:g10:orders-of-magnitude:unit}{ثانية}، وهو جدير بالتذكر قبل النزول من الرصيف.
والصوت في الهواء: $\qty{340}{m/s} \times 3.6 = \qty{1224}{km/h}$.
\end{example}

\begin{example}[نزهة بفاصل راحة]\label{ex:g10:relative-motion:hike}
يقطع متنزّه \qty{12}{km} في \qty{2}{h} من المشي، ثم يستريح
\qty{1}{h}، ثم يقطع \qty{6}{km} في \qty{1.5}{h} الأخيرة. \omterm{def:g10:relative-motion:speed}{والسرعة
المتوسطة} على النزهة كلها:
\[
v = \frac{d}{t} = \frac{12 + 6}{2 + 1 + 1.5}
= \frac{\qty{18}{km}}{\qty{4.5}{h}} = \qty{4.0}{km/h}.
\]
فالمتوسط يحسب \emph{كل} الزمن المنقضي، بما فيه فترات الراحة --- وهو
عمومًا ليس متوسط سرعات المراحل المنفصلة.
\end{example}

\section{اختيار مرجع: الأرضي والمركزي الأرضي والمركزي الشمسي}

بما أن كل \omterm{def:g10:relative-motion:frame}{مرجع} مشروع، فإننا نختار \omterm{def:g10:relative-motion:frame}{المرجع} الذي تبدو فيه الحركة
التي تهمّنا أبسط ما تكون. وثلاثة اختيارات تغطي معظم فيزياء
هذا الكتاب.

\begin{definition}[المراجع الثلاثة المعيارية]\label{def:g10:relative-motion:frames}
\begin{itemize}
  \item \emph{المرجع الأرضي}\index{مرجع أرضي} (أو
        \emph{مرجع المختبر}\index{مرجع المختبر}) مرتبط
        بسطح الأرض حيث تجري التجربة. وهو \omterm{def:g10:relative-motion:frame}{المرجع}
        المناسب للحركة اليومية: الأجسام الساقطة، والمركبات،
        والرياضة، وطاولات المختبر.
  \item \emph{المرجع المركزي الأرضي}\index{مرجع مركزي أرضي} مركزه
        الأرض لكنه \emph{لا} يدور معها (فاتجاهاته
        تشير إلى نجوم بعيدة). وهو \omterm{def:g10:relative-motion:frame}{المرجع} المناسب للقمر
        وللتوابع الاصطناعية، وهي تدور حول الأرض كلها
        وتتجاهل دورانها اليومي.
  \item \emph{المرجع المركزي الشمسي}\index{مرجع مركزي شمسي}
        مركزه الشمس، واتجاهاته أيضًا مثبَّتة على نجوم
        بعيدة. وهو \omterm{def:g10:relative-motion:frame}{المرجع} المناسب للكواكب: فيتبع كل منها عندئذٍ
        مدارًا بسيطًا شبه دائري.
\end{itemize}
\end{definition}

\begin{omfigure}
\begin{tikzpicture}[scale=1.0, font=\small]
  \fill[omProp] (0,0) circle (0.26);
  \node[below=3pt] at (0,-0.26) {الشمس};
  \draw[dashed, gray] (0,0) circle (2.5);
  \draw[-Stealth, gray] (2.5,0) arc (0:22:2.5);
  \fill[omDef] (2.5,0) circle (0.12);
  \node[below right] at (2.55,-0.10) {الأرض};
  \draw[dashed, gray] (2.5,0) circle (0.75);
  \draw[-Stealth, gray] (3.25,0) arc (0:55:0.75);
  \fill[gray] (2.5,0.75) circle (0.06);
  \node[above] at (2.5,0.80) {القمر};
  \draw[omThm] (2.38,0.12) -- (1.2,1.35);
  \node[omThm, left, align=right] at (1.25,1.45)
    {المرجع الأرضي:\\نقطة تدور\\مع السطح};
  \node[omDef, align=left, anchor=west] at (3.6,0.1)
    {المرجع المركزي الأرضي:\\مركزه الأرض\\(القمر، التوابع)};
  \node[omProp, align=left, anchor=west] at (-1.4,-1.6)
    {المرجع المركزي الشمسي: مركزه الشمس (الكواكب)};
\end{tikzpicture}

{\small ثلاثة مراجع متداخلة: الأرضي يدور مع الأرض،
والمركزي الأرضي يركب مع الأرض دون أن يدور، والمركزي
الشمسي يشاهد الرقصة كلها من الشمس.}
\end{omfigure}

\begin{method}[اختيار المرجع]\label{met:g10:relative-motion:choose}
سمِّ الأجسام التي تريد وصف حركتها، ثم اختر \omterm{def:g10:relative-motion:frame}{المرجع}
الذي تكون فيه تلك الحركة أبسط ما تكون:
\begin{enumerate}
  \item حركة قرب سطح الأرض تدوم دقائق --- \omterm{def:g10:relative-motion:frames}{المرجع
        الأرضي}؛
  \item مدارات \emph{حول} الأرض (القمر، التوابع) --- \omterm{def:g10:relative-motion:frames}{المرجع
        المركزي الأرضي}؛
  \item مدارات حول الشمس (الكواكب، المذنبات) --- \omterm{def:g10:relative-motion:frames}{المرجع المركزي الشمسي}.
\end{enumerate}
واذكر اختيارك دائمًا قبل وصف حركة: \omterm{def:g10:relative-motion:trajectory}{فالمسار} أو
السرعة المذكورة بلا مرجعها عديمة المعنى
(\cref{prop:g10:relative-motion:relativity}).
\end{method}

\begin{remark}[من المحقّ، بطليموس أم كوبرنيكوس؟]\label{rem:g10:relative-motion:history}
``تشرق الشمس من الشرق'' و``تدور الأرض نحو الشرق'' هما
الرصد نفسه مبلَّغًا في مرجعين --- الأول في \omterm{def:g10:relative-motion:frames}{المرجع
الأرضي}، حيث تكنس الشمس السماء، والثاني في \omterm{def:g10:relative-motion:frames}{المرجع
المركزي الأرضي}، حيث يُحمل الراصد حول محور الأرض. وقد ظلت
المراجع نفسها ميدان المعركة قرونًا: فوصف بطليموس
المتمركز حول الأرض تعقّب السماء بدقة لكنه احتاج حلقات
فوق حلقات لتتبّع الكواكب، بينما جعل الوصفُ المتمركز حول الشمس
الذي أحياه كوبرنيكوس سنة 1543 كلَّ كوكب شبه دائرة بسيطة،
وفسّر حلقات المريخ الخلفية الغريبة أثرًا مرجعيًّا لا غير ---
أي تجاوز الأرض كوكبًا أبطأ. والحكم الحديث ليس أن مرجعًا
صادق والآخر كاذب، بل أنهما ليسا متساويين في
\emph{الملاءمة}؛ أما معيار أعمق لتفضيل بعض المراجع
(المراجع العُطالية) فيأتي مع مبدأ العطالة في
\cref{ch:g10:inertia}.
\end{remark}

\section{تركيب الحركات على مستقيم}

ما الذي يصل \qty{4}{km/h} للماشي (في \omterm{def:g10:relative-motion:frame}{مرجع} القطار) بالسرعة
\qty{304}{km/h} (في \omterm{def:g10:relative-motion:frames}{المرجع الأرضي})؟ على مستقيم واحد، الجواب جمع
بسيط.

\begin{proposition}[تركيب السرعات في بُعد واحد]\label{prop:g10:relative-motion:compose}
لنفترض جسمًا يتحرك على مستقيم فوق حامل متحرك (سير متحرك،
أو قطار، أو ماء جارٍ)، والحامل نفسه ينزلق على المستقيم نفسه
بالنسبة إلى الأرض. وباختيار جهة موجبة على المستقيم
وعدّ السرعات بإشارة،
\[
v_{\text{الجسم/الأرض}} \;=\; v_{\text{الجسم/الحامل}} \;+\;
v_{\text{الحامل/الأرض}} .
\]
وبالكلام: السرعتان في الجهة نفسها تتجمعان؛ والسرعتان في جهتين
متعاكستين تتطارحان.
\end{proposition}

\begin{proof}
خلال \omterm{def:g10:orders-of-magnitude:unit}{ثانية} واحدة، يحمل الحاملُ الجسمَ $v_{\text{الحامل/الأرض}}$
مترًا على المستقيم، ويتقدم الجسم $v_{\text{الجسم/الحامل}}$
مترًا على الحامل. والانتقالان يقعان على المستقيم نفسه،
ومن ثم يكون انتقال الجسم الكلي على الأرض في تلك \omterm{def:g10:orders-of-magnitude:unit}{الثانية} مجموعَهما
الجبري --- وهو سرعته بالنسبة إلى الأرض.
\end{proof}

\begin{example}[السير المتحرك]\label{ex:g10:relative-motion:walkway}
يتحرك سير في مطار بسرعة \qty{1.0}{m/s}. ويمشي مسافر عليه بسرعة
\qty{1.5}{m/s} بالنسبة إلى السير، في جهة السير: فتكون سرعته
بالنسبة إلى الأرض $1.5 + 1.0 = \qty{2.5}{m/s}$. وإذا مشى في الجهة
الخاطئة بالخطى نفسها: $1.5 - 1.0 = \qty{0.5}{m/s}$، فيظل متقدمًا (ببطء)
عكس السير. وإذا وقف ساكنًا على السير:
\qty{1.0}{m/s} بالنسبة إلى الأرض، و\qty{0}{m/s} بالنسبة إلى
السير --- أي ساكنًا ومتحركًا في آن، في مرجعين
مختلفين.
\end{example}

\begin{omfigure}
\begin{tikzpicture}[font=\small]
  \draw[fill=gray!12] (0,0) rectangle (7,0.3);
  \node[gray] at (6.3,0.15) {السير};
  \draw (2.5,0.3) -- (2.5,1.05);
  \draw (2.5,0.75) -- (2.2,0.45) (2.5,0.75) -- (2.8,0.45);
  \fill (2.5,1.2) circle (2.5pt);
  \draw[-Stealth, omDef, thick] (0.4,-0.45) -- (1.9,-0.45)
    node[midway, below] {السير/الأرض: \qty{1.0}{m/s}};
  \draw[-Stealth, omThm, thick] (2.5,1.6) -- (4.75,1.6)
    node[midway, above] {الماشي/السير: \qty{1.5}{m/s}};
  \draw[-Stealth, omMeth, very thick] (2.5,2.3) -- (6.25,2.3)
    node[midway, above] {الماشي/الأرض: \qty{2.5}{m/s}};
\end{tikzpicture}

{\small التركيب في بُعد واحد: السهام مرسومة بالمقياس،
وسهم \omterm{def:g10:relative-motion:frames}{المرجع الأرضي} هو مجموع السهمين الآخرين.}
\end{omfigure}

\begin{remark}[اتجاهان في آن]\label{rem:g10:relative-motion:twodim}
حين لا تكون الحركتان \emph{على} المستقيم نفسه --- كقارب يصوّب
عبورًا مستقيمًا لنهر بينما يدفعه التيار مع الجريان --- تجري كل
حركة مجراها مستقلة، وتتركّب السرعتان
سهمين موضوعين رأسًا إلى ذيل، تمامًا كمتجهات مجلد
الرياضيات. ويظل القارب يبلغ الضفة البعيدة في الزمن الذي تحدّده
سرعة عبوره وحدها، لكنه يهبط مع الجريان بعد نقطة تصويبه،
ويكون مساره في \omterm{def:g10:relative-motion:frames}{المرجع الأرضي} مستقيمًا مائلًا. أما المعالجة المتجهية
الكاملة فتأتي في فصل لاحق؛ وهنا لا نقرأ إلا الصورة.
\end{remark}

\begin{omfigure}
\begin{tikzpicture}[scale=0.95, font=\small]
  \fill[omDef!10] (0,0) rectangle (6.4,2.4);
  \draw (0,0) -- (6.4,0) (0,2.4) -- (6.4,2.4);
  \draw[-Stealth, gray] (4.8,0.35) -- (5.6,0.35);
  \node[gray] at (5.2,0.62) {التيار};
  \draw[dashed, omMeth] (1.5,0) -- (3.3,2.4);
  \draw[-Stealth, omThm, thick] (1.5,0.6) -- (1.5,2.0)
    node[left, pos=0.7] {القارب/الماء};
  \draw[-Stealth, omDef, thick] (1.5,0.6) -- (2.55,0.6)
    node[below, pos=0.8] {الماء/الأرض};
  \draw[-Stealth, omMeth, very thick] (1.5,0.6) -- (2.55,2.0)
    node[right, pos=0.75] {القارب/الأرض};
  \fill (1.5,0.6) circle (1.8pt);
\end{tikzpicture}

{\small قارب يتجه عبورًا مستقيمًا لنهر: ففي \omterm{def:g10:relative-motion:frame}{مرجع} الماء
يعبر مستقيمًا؛ وفي \omterm{def:g10:relative-motion:frames}{المرجع الأرضي} يتبع الخط المتقطع المائل،
منجرفًا مع الجريان وهو يعبر.}
\end{omfigure}

\section{تمارين}

\begin{exercise}[$\star$]\label{exo:g10:relative-motion:1}
مسافر نائم في قطار يسير بسرعة \qty{100}{km/h}. قل
أساكن هو أم متحرك، وبأي سرعة، بالنسبة إلى: العربة؛ والأرض؛
وشجرة بجانب السكة؛ وقطار قادم يسير بسرعة \qty{100}{km/h} في الاتجاه
المعاكس. ثم صِف حركة الشجرة في \omterm{def:g10:relative-motion:frame}{مرجع} قطار المسافر.
\end{exercise}

\begin{exercise}[$\star$]\label{exo:g10:relative-motion:2}
حوّل: \qty{72}{km/h} و\qty{108}{km/h} إلى \unit{m/s}؛
و\qty{25}{m/s} و\qty{340}{m/s} إلى \unit{km/h}. وتحقّق من كل جواب
بقاعدة السلامة في \cref{met:g10:relative-motion:convert}.
\end{exercise}

\begin{exercise}[$\star$]\label{exo:g10:relative-motion:3}
يقطع راكب دراجة \qty{36}{km} في \qty{1}{h}~\qty{30}{min}؛ ويقطع عدّاء
\qty{100}{m} في \qty{10.0}{s}. احسب السرعتين المتوسطتين \omterm{def:g10:orders-of-magnitude:unit}{بوحدة}
\unit{m/s} وبوحدة \unit{km/h}. فمن الأسرع، وهل يفاجئك
الجواب؟
\end{exercise}

\begin{exercise}[$\star$]\label{exo:g10:relative-motion:4}
يتحرك سير بسرعة \qty{0.9}{m/s}. ويمشي مسافر بسرعة
\qty{1.4}{m/s} بالنسبة إلى السير. أعطِ سرعة المسافر بالنسبة
إلى الأرض حين يمشي في جهة السير، وعكسها، وحين
يقف ساكنًا على السير.
\end{exercise}

\begin{exercise}[$\star$]\label{exo:g10:relative-motion:5}
لكل حركة مما يلي، اختر \omterm{def:g10:relative-motion:frame}{المرجع} الأنسب من بين الأرضي
والمركزي الأرضي والمركزي الشمسي، وبرّر في جملة واحدة: قلم يسقط
داخل قطار متحرك (وزيادة: ثمة \omterm{def:g10:relative-motion:frame}{مرجع} أفضل بعدُ)؛ \omterm{def:g10:universal-gravitation:satellite}{ومدار}
محطة الفضاء الدولية؛ وسنة مريخية واحدة؛ وسباق
عدو \qty{100}{m}.
\end{exercise}

\begin{exercise}[$\star\star$]\label{exo:g10:relative-motion:6}
يسير قطار بسرعة \qty{108}{km/h}. ويمشي مسافر نحو المقدمة بسرعة
\qty{1.5}{m/s} بالنسبة إلى القطار.
\begin{enumerate}
  \item حوّل سرعة القطار إلى \unit{m/s}، ثم أعطِ
        سرعة المسافر بالنسبة إلى الأرض، ماشيًا نحو
        المقدمة ونحو المؤخرة.
  \item طول العربة \qty{60}{m}. فكم يستغرق المشي إلى
        المقدمة، وكم يقطع المسافر بالنسبة إلى الأرض خلاله؟
\end{enumerate}
\end{exercise}

\begin{exercise}[$\star\star$]\label{exo:g10:relative-motion:7}
سرعة قارب في ماء ساكن \qty{5.0}{m/s}؛ ويجري نهر بسرعة
\qty{2.0}{m/s}. ويسير القارب \qty{420}{m} مع الجريان، ثم يستدير
ويعود.
\begin{enumerate}
  \item احسب سرعة القارب بالنسبة إلى الضفة في كل مرحلة،
        ومدة كل مرحلة.
  \item احسب \omterm{def:g10:relative-motion:speed}{السرعة المتوسطة} في الرحلة كلها. ولماذا هي أقلّ
        من \qty{5.0}{m/s}، مع أن التيار يساعد في مرحلة
        بمقدار ما يعوق في الأخرى بالضبط؟
\end{enumerate}
\end{exercise}

\begin{exercise}[$\star\star$]\label{exo:g10:relative-motion:8}
يتجاوز قطاران أحدهما الآخر على سكتين متوازيتين: القطار A (طوله
\qty{180}{m}) بسرعة \qty{25}{m/s}، والقطار B (طوله \qty{90}{m}) بسرعة \qty{20}{m/s}.
\begin{enumerate}
  \item يسيران في جهتين متعاكستين. فما سرعة B في
        \omterm{def:g10:relative-motion:frame}{مرجع} A، وكم يدوم التجاوز الكامل
        (من التقاء المقدمتين إلى انفصال المؤخرتين)؟
  \item والسؤالان نفسهما إذا سارا في الجهة نفسها.
\end{enumerate}
\end{exercise}

\begin{exercise}[$\star\star$]\label{exo:g10:relative-motion:9}
تسير دراجة على طريق مستقيم. صِف \omterm{def:g10:relative-motion:trajectory}{مسار}
صمام الإطار في: \omterm{def:g10:relative-motion:frame}{مرجع} راكب الدراجة؛ \omterm{def:g10:relative-motion:frames}{والمرجع الأرضي}؛ \omterm{def:g10:relative-motion:frame}{ومرجع}
الصمام نفسه. وصِف كذلك \omterm{def:g10:relative-motion:trajectory}{مسار} محور العجلة في
\omterm{def:g10:relative-motion:frames}{المرجع الأرضي}. ولا حساب --- إنما سمِّ كل منحنٍ.
\end{exercise}

\begin{exercise}[$\star\star$]\label{exo:g10:relative-motion:10}
طول خط استواء الأرض نحو \qty{40000}{km}، وتدور الأرض
دورة في \qty{24}{h}.
\begin{enumerate}
  \item احسب سرعة نقطة على خط الاستواء في \omterm{def:g10:relative-motion:frames}{المرجع المركزي
        الأرضي}، \omterm{def:g10:orders-of-magnitude:unit}{بوحدة} \unit{km/h} وبوحدة \unit{m/s}.
  \item اشرح في جملة واحدة لكل: لماذا لا نحسّ هذه السرعة،
        ولماذا تبلّغ عبارتا ``تشرق الشمس'' و``تدور الأرض'' الرصد
        نفسه.
\end{enumerate}
\end{exercise}

\begin{exercise}[$\star\star$]\label{exo:g10:relative-motion:11}
بالوقوف ساكنًا على سلّم متحرك، يستغرق الصعود \qty{60}{s}.
والصعود مشيًا على السلّم نفسه مطفأً يستغرق \qty{90}{s}. فكم
يستغرق الصعود مشيًا على السلّم المتحرك؟ (اشتغل بالسرعات لا
بالأزمنة.)
\end{exercise}

\begin{exercise}[$\star\star\star$]\label{exo:g10:relative-motion:12}
يقطع سائق المسافة نفسها مرتين: ذهابًا بسرعة \qty{60}{km/h}،
وإيابًا بسرعة \qty{90}{km/h}.
\begin{enumerate}
  \item بيّن أن \omterm{def:g10:relative-motion:speed}{السرعة المتوسطة} في الرحلة كلها
        \qty{72}{km/h} --- لا \qty{75}{km/h}.
  \item اشرح أين يخطئ الجواب الساذج، وقل أيّ المرحلتين
        تهيمن على المتوسط.
\end{enumerate}
\end{exercise}

\begin{exercise}[$\star\star\star$]\label{exo:g10:relative-motion:13}
ينطلق راكبا دراجتين تفصل بينهما \qty{35}{km} ويركبان أحدهما نحو الآخر، أحدهما بسرعة
\qty{15}{km/h}، والآخر بسرعة \qty{20}{km/h}.
\begin{enumerate}
  \item صِف حركة راكب الدراجة الثاني في \omterm{def:g10:relative-motion:frame}{مرجع}
        الأول، واستنتج متى يلتقيان.
  \item كم قطع كل منهما عند اللقاء؟
  \item وإذا ركبا بدل ذلك في الجهة نفسها (والأسرع
        خلف)، فكم تدوم المطاردة؟
\end{enumerate}
\end{exercise}

\begin{exercise}[$\star\star\star$]\label{exo:g10:relative-motion:14}
يعبر سابح نهرًا عرضه \qty{60}{m}، متجهًا عبورًا مستقيمًا بسرعة
\qty{1.2}{m/s} بالنسبة إلى الماء؛ ويجري التيار بسرعة
\qty{0.9}{m/s}.
\begin{enumerate}
  \item اشرح لماذا يحدّد زمنَ العبور سرعةُ عبور السابح
        وحدها، واحسبه.
  \item كم يهبط السابح مع الجريان؟
  \item باستعمال مثلث قائم، جِد سرعة السابح في \omterm{def:g10:relative-motion:frames}{المرجع
        الأرضي} وطول \omterm{def:g10:relative-motion:trajectory}{المسار} المسبوح فعلًا على الأرض.
\end{enumerate}
\end{exercise}

\begin{exercise}[$\star\star\star$]\label{exo:g10:relative-motion:15}
تدور الأرض حول الشمس على شبه دائرة نصف قطرها
\qty{1.50e11}{m}، في سنة واحدة ($\num{3.16e7}$ \omterm{def:g10:orders-of-magnitude:unit}{ثانية}).
\begin{enumerate}
  \item احسب سرعة الأرض في \omterm{def:g10:relative-motion:frames}{المرجع المركزي الشمسي}، \omterm{def:g10:orders-of-magnitude:unit}{بوحدة}
        \unit{km/s}.
  \item قارنها بسرعة دوران خط الاستواء التي وجدتها في
        \cref{exo:g10:relative-motion:10}.
  \item تبدو المريخ، منظورًا إليه من الأرض، وكأنه يتوقف أحيانًا وينجرف
        إلى الوراء بين النجوم بضعة أسابيع. اشرح في جملتين،
        بمفردات هذا الفصل، لماذا يذيب \omterm{def:g10:relative-motion:frames}{المرجع المركزي
        الشمسي} هذا اللغز.
\end{enumerate}
\end{exercise}

\section{مسألة: السير والنهر والتابع}

\begin{problem}[{مسألة نهاية الأسبوع --- السير المتحرك والنهر
والتابع: من يتحرك، وبالنسبة إلى ماذا، وتقرير السرعة النهائي
للمسافر القاعد}]
\label{pb:g10:relative-motion:1}
مطار، وعبور بعبّارة، ونقطة مضيئة تنزلق عبر سماء الغسق:
ثلاثة مشاهد يومية، وسؤال واحد يُطرح ثلاث مرات --- من
يتحرك، وبالنسبة إلى ماذا؟ ويختار كل جزء مرجعه، ويحسب،
ويرقب الجواب يتغير مع وجهة النظر؛ وتحسب الخاتمة
كم سرعة \emph{حركتك} أنت الآن، وأنت جالس ساكنًا تمامًا.

\medskip
\noindent\textbf{الجزء I --- سباق السير المتحرك.}
سير في مطار طوله \qty{90}{m} ويتحرك بسرعة \qty{1.0}{m/s}. ويمشي
مسافر بسرعة \qty{1.5}{m/s} (بالنسبة إلى ما تحت قدميه).
\begin{enumerate}
  \item كم يستغرق السير لينقل مسافرًا واقفًا
        من طرف إلى طرف؟
  \item وكم يستغرق المسافر إذا مشى على السير
        المتحرك، في جهته؟
  \item وكم يستغرق إذا مشى بجانب السير، على الأرضية الثابتة؟
  \item يمشي مازح في الجهة الخاطئة على السير، بسرعة \qty{1.5}{m/s}
        بالنسبة إليه. فما سرعته بالنسبة إلى الأرض، وكم يستغرق ليعود
        من الطرف البعيد إلى المدخل؟
  \item يركض طفل في الجهة الخاطئة بسرعة \qty{1.0}{m/s} بالضبط بالنسبة
        إلى السير. صِف حركة الطفل في \omterm{def:g10:relative-motion:frames}{المرجع الأرضي}،
        وسمِّ آلة الصالة الرياضية التي يعيد اختراعها.
  \item لخّص الأسئلة 1--5 في جدول من عمودين: سرعة كل متحرك
        في \omterm{def:g10:relative-motion:frame}{مرجع} السير وفي \omterm{def:g10:relative-motion:frames}{المرجع الأرضي}. فأيّ العمودين
        تحسّه أرجل المسافرين؟
\end{enumerate}

\medskip
\noindent\textbf{الجزء II --- النهر.}
تعبر عبّارة نهرًا عرضه \qty{120}{m}. وسرعتها بالنسبة إلى
الماء \qty{2.0}{m/s}؛ ويجري التيار بسرعة \qty{1.5}{m/s}. ويصوّب الربّان
أولًا مستقيمًا نحو الضفة المقابلة.
\begin{enumerate}[resume]
  \item كم يستغرق العبور؟
  \item كم تهبط العبّارة مع الجريان؟
  \item باستعمال مثلث قائم، احسب سرعة العبّارة في \omterm{def:g10:relative-motion:frames}{المرجع
        الأرضي} وطول مسارها في \omterm{def:g10:relative-motion:frames}{المرجع الأرضي}. (وسيظهر
        ثلاثي أعداد مألوف.)
  \item صِف \omterm{def:g10:relative-motion:trajectory}{مسار} العبّارة في \omterm{def:g10:relative-motion:frame}{مرجع} الماء وفي
        \omterm{def:g10:relative-motion:frames}{المرجع الأرضي}. فما المشترك بينهما، وما المختلف؟
  \item ينطلق سابح سرعته بالنسبة إلى الماء \qty{1.0}{m/s}
        من الضفة نفسها، مصوّبًا عبورًا مستقيمًا مثل
        العبّارة. فكم يستغرق العبور، وكم يهبط مع الجريان؟
        احسب سرعة السابح في \omterm{def:g10:relative-motion:frames}{المرجع الأرضي} (بمثلث قائم مرة
        أخرى). ولماذا ينجرف السابح الأبطأ أبعد من العبّارة؟
  \item يمرّ جذع خشب طافيًا بجانب العبّارة. أعطِ سرعته في \omterm{def:g10:relative-motion:frames}{المرجع الأرضي}
        وفي \omterm{def:g10:relative-motion:frame}{مرجع} الماء. ففي أي \omterm{def:g10:relative-motion:frame}{مرجع} يكون الجذع ``متحركًا''؟
        علّق.
\end{enumerate}

\medskip
\noindent\textbf{الجزء III --- مرور \omterm{def:g10:universal-gravitation:satellite}{التابع}.}
عند الغسق، تعبر محطة الفضاء الدولية السماء في بضع
دقائق. وهي تدور على ارتفاع \qty{400}{km}؛ ونصف قطر الأرض
\qty{6.37e6}{m}، وتستغرق الدورة الواحدة \qty{5.56e3}{s} (أي نحو
\qty{93}{min}).
\begin{enumerate}[resume]
  \item في أي \omterm{def:g10:relative-motion:frame}{مرجع} معياري تكون حركة المحطة أبسط؟
        احسب سرعتها في ذلك \omterm{def:g10:relative-motion:frame}{المرجع}، \omterm{def:g10:orders-of-magnitude:unit}{بوحدة} \unit{km/s}.
  \item اشرح في جملتين لماذا يصف \omterm{def:g10:relative-motion:frames}{المرجع الأرضي}
        المحطة وصفًا عسيرًا --- ماذا يحدث لأثرها على الأرض
        بين دورة والتي تليها؟
  \item قارن سرعة المحطة المدارية بسرعة دوران
        خط الاستواء ($\approx \qty{0.46}{km/s}$): بأي عامل تكون
        المحطة أسرع؟
  \item يدور \omterm{def:g10:universal-gravitation:satellite}{تابع} ثابت بالنسبة إلى الأرض على نصف قطر \qty{4.22e7}{m} في
        يوم واحد بالضبط ($\qty{8.64e4}{s}$). احسب سرعته في
        \omterm{def:g10:relative-motion:frames}{المرجع المركزي الأرضي}، ثم صِف حركته في \omterm{def:g10:relative-motion:frames}{المرجع الأرضي}.
        ووفّق بين الجوابين في جملة واحدة.
  \item في \omterm{def:g10:relative-motion:frames}{المرجع المركزي الأرضي} يرسم القمر شبه دائرة حول
        الأرض في نحو $27$ يومًا. صِف بالكلام مساره في
        \omterm{def:g10:relative-motion:frames}{المرجع المركزي الشمسي} على مدى بضعة أشهر.
\end{enumerate}

\medskip
\noindent\textbf{الجزء IV --- من يتحرك؟}
\begin{enumerate}[resume]
  \item ``تشرق الشمس من الشرق.'' ``تدور الأرض نحو
        الشرق.'' اذكر \omterm{def:g10:relative-motion:frame}{المرجع} الذي تصف فيه كل جملة
        الرصد، واشرح لماذا ليست أيّ منهما خاطئة.
  \item في فقرة واحدة رصينة: ما الذي أحسنه \omterm{def:g10:relative-motion:frame}{مرجع} بطليموس، وما
        الذي بسّطه \omterm{def:g10:relative-motion:frame}{مرجع} كوبرنيكوس (اذكر حلقات المريخ
        الخلفية)، ولماذا يدور الحكم الحديث حول \emph{الملاءمة}
        لا حول كون \omterm{def:g10:relative-motion:frame}{مرجع} صادقًا؟
  \item تقرير سرعة المسافر القاعد. أنت جالس ساكنًا في
        كرسي قرب خط الاستواء. أعطِ سرعتك في \omterm{def:g10:relative-motion:frames}{المرجع الأرضي}،
        وفي \omterm{def:g10:relative-motion:frames}{المرجع المركزي الأرضي} ($\approx \qty{0.46}{km/s}$)،
        وفي \omterm{def:g10:relative-motion:frames}{المرجع المركزي الشمسي} ($\approx \qty{30}{km/s}$). ثم احسب
        المسافة التي تقطعها في \omterm{def:g10:relative-motion:frames}{المرجع المركزي الشمسي} خلال
        فيلم مدته \qty{90}{min}، واختم المسألة بحكمة الفصل
        في جملة واحدة.
\end{enumerate}
\end{problem}
